\documentclass[preprint,12pt]{elsarticle}

% --- Packages ---
\usepackage{amsmath,amssymb}
\usepackage{booktabs}
\usepackage{graphicx}
\usepackage{hyperref}
\usepackage{xcolor}
\usepackage{lineno}
\usepackage{multirow}
\usepackage{siunitx}

% Enable line numbers for review
\modulolinenumbers[5]

% Figures are in figures/ subdirectory
\graphicspath{{figures/}}

\journal{Acta Astronautica}

\begin{document}

\begin{frontmatter}

\title{Microgravity Metallurgy at Industrial Scale: Scaling Laws, Architecture Trade-offs, and Technology Maturation Pathways for In-Space Metal Processing}

\author[dyson]{Thijs Hakkenberg\corref{cor1}}
\cortext[cor1]{Corresponding author.}
\ead{thijs@hakkenberg.com}
\fntext[fn1]{This research was conducted using an AI-assisted multi-model consensus methodology described in \cite{hakkenberg2026consensus}. Technical proposals from Claude, Gemini, and GPT were synthesized and validated against published literature.}

\address[dyson]{Project Dyson, Open Research Initiative\fnref{fn1}}

\begin{abstract}
Large-scale space infrastructure programs require tens of thousands of tonnes of processed metals annually, yet demonstrated microgravity metallurgy is limited to gram-scale experiments aboard the International Space Station. This paper examines whether the six-to-eight orders-of-magnitude gap between demonstrated capability and required throughput can be bridged, and if so, through what architectural pathway. We review the state of the art in microgravity materials processing, analyze the gravity sensitivity of each unit operation in the metal processing chain, and evaluate a hybrid multi-gravity-zone architecture against pure microgravity and full artificial gravity alternatives.

Our central finding is that the processing chain should be decomposed by gravity sensitivity: primary smelting and slag separation in a rotating module at 0.05--0.15$g$ where buoyancy-driven phase separation functions, and downstream refining and thin-film deposition in microgravity where containerless processing and suppressed convection offer demonstrated quality advantages. This hybrid architecture fits within a 400,000--500,000~kg station mass envelope, adding only 10,000--15,000~kg net penalty over a pure microgravity design after accounting for eliminated electromagnetic confinement hardware. We define quantitative go/no-go criteria for a Gate~1 technology decision at month~36, identify the critical experiments needed in the 0.01--0.2$g$ regime where almost no data exists, and propose a phased \$550--810M research program to retire the highest-risk unknowns.
\end{abstract}

\begin{keyword}
microgravity metallurgy \sep in-space manufacturing \sep artificial gravity \sep zone refining \sep ISRU \sep space station design
\end{keyword}

\end{frontmatter}

\linenumbers

%% ============================================================
\section{Introduction}
\label{sec:introduction}
%% ============================================================

Any program to construct large-scale space infrastructure from in-situ resources must solve the metallurgy problem: how to convert raw asteroidal or lunar feedstock into structural metals, solar-grade silicon, and functional alloys at production rates measured in tens of thousands of tonnes per year. This processing challenge is widely acknowledged as a project-ending risk for ambitious space manufacturing concepts, yet quantitative assessments of the scaling pathway are rare in the literature.

The challenge is not merely one of scaling up existing processes. Terrestrial metallurgy relies fundamentally on gravity for slag--metal separation, convective heat transfer in furnaces, and melt containment in refractory-lined vessels. In microgravity, all three mechanisms are absent. Electromagnetic levitation can contain small molten metal samples---the ISS Electromagnetic Levitator (EML) processes specimens of order 1~cm$^3$---but the power requirements for electromagnetic containment scale unfavorably with sample mass, and thermal management without buoyancy-driven convection becomes increasingly problematic at larger scales.

This paper addresses three questions. First, what is the actual state of the art in microgravity metallurgy, and what can be reliably extrapolated from it (\S\ref{sec:state-of-art})? Second, which unit operations in the metal processing chain are genuinely gravity-dependent, which benefit from microgravity, and what is the minimum gravity level needed for each (\S\ref{sec:gravity-sensitivity})? Third, given the answers to these questions, what station architecture and research program can most efficiently retire the risk (\S\ref{sec:architecture}, \S\ref{sec:gate1})?

Our analysis draws on a systematic literature review of microgravity materials science experiments, multi-model AI deliberations that synthesized proposals from three independent AI systems, and parametric scaling analysis of key process constraints. The deliberation on this question---conducted under the multi-model consensus methodology described in \cite{hakkenberg2026consensus}---converged on a hybrid multi-gravity-zone architecture as the clear preferred approach, a finding we develop and quantify here.

%% ============================================================
\section{State of the Art in Microgravity Materials Processing}
\label{sec:state-of-art}
%% ============================================================

\subsection{ISS Electromagnetic Levitator experiments}

The ISS EML, operated by DLR under ESA's Materials Science Laboratory program, represents the most capable orbital facility for containerless metal processing. The EML uses electromagnetic coils to position, heat, and melt metallic samples of approximately 6--8~mm diameter (0.5--2~g mass) under ultra-high vacuum or controlled gas atmospheres. As of 2025, the EML has processed over 300 samples spanning nickel-based superalloys, steel, titanium alloys, and semiconductor materials \cite{ratke1995}.

The scientific value of these experiments is considerable---they have provided unique data on nucleation, undercooling, and solidification in the absence of buoyancy-driven convection. However, from an industrial scaling perspective, two limitations are critical:

\begin{enumerate}
    \item \textbf{Sample mass:} The EML is designed for samples below 3~g. Scaling electromagnetic containment to even 1~kg samples would require approximately $10^3$~times the coil current or a fundamentally different coil geometry, with concomitant power increases.
    \item \textbf{Processing mode:} The EML is a batch research instrument with manual sample loading. Continuous or semi-continuous processing---essential for industrial throughput---has never been demonstrated in orbit.
\end{enumerate}

\subsection{Gravity-independent additive manufacturing}

D'Angelo et al.\ \cite{dangelo2021} demonstrated a powder-based additive manufacturing process designed to function independently of gravitational orientation. Using controlled powder delivery mechanisms that do not rely on gravity-fed hoppers, they produced metal parts in parabolic flight campaigns with quality comparable to 1$g$ reference samples. This represents a genuine advance for in-space manufacturing of finished components, but it addresses fabrication rather than the upstream bulk smelting and refining that constitute the volume bottleneck.

\subsection{Ultrasound-assisted grain refinement}

Todaro et al.\ \cite{todaro2020} showed that high-intensity ultrasound applied during solidification in additive manufacturing produces fine, equiaxed grain structures comparable to those achieved through gravity-driven convection. This finding is significant because it demonstrates that at least one mechanism for controlling microstructure in microgravity exists and is practical. However, the demonstrated scale remains below 1~kg, and the applicability to bulk smelting rather than AM deposition is unproven.

\subsection{Microgravity electrolysis}

Akay, Romero-Calvo et al.\ \cite{akay2025} demonstrated that a commercial neodymium magnet can enhance water electrolysis efficiency in microgravity by up to 240\%, exploiting diamagnetic buoyancy and magnetohydrodynamic forces for passive gas--liquid phase separation. This breakthrough, validated in drop-tower experiments \cite{romerocalo2022}, eliminates the primary barrier to microgravity electrolysis (mechanical separation of product gases from electrolyte). Scaling from laboratory prototypes to the 500--750~kW electrolysis systems required for industrial hydrogen and oxygen production remains the principal open challenge.

\subsection{The scaling gap}

Table~\ref{tab:scaling-gap} summarizes the gap between demonstrated capability and required production rates for each unit operation.

\begin{table}[ht]
\centering
\caption{Scaling gap between demonstrated microgravity metallurgy and required industrial throughput.}
\label{tab:scaling-gap}
\begin{tabular}{llll}
\toprule
\textbf{Operation} & \textbf{Demonstrated} & \textbf{Required} & \textbf{Gap} \\
\midrule
Smelting/melting & $\sim$2~g (EML) & 5--20~tonnes/hr & $\sim 10^{7}\times$ \\
Zone refining & Not demonstrated & 100~kg/batch & $> 10^{5}\times$ \\
Electrolysis & $\sim$10~mL/hr & 500--750~kW & $\sim 10^{4}\times$ \\
Additive manufacturing & $\sim$0.5~kg parts & 100~kg parts & $\sim 10^{2}\times$ \\
Thin-film deposition & $\sim$0.1~m$^2$ & 10--100~m$^2$ & $\sim 10^{2}\times$ \\
\bottomrule
\end{tabular}
\end{table}

The gap is not uniform across operations. Additive manufacturing and thin-film deposition face scale-up challenges of 2--3 orders of magnitude that are plausibly bridgeable through engineering development. Smelting and zone refining face 5--7 orders of magnitude, a gap that cannot be closed by incremental improvements to existing microgravity techniques.

%% ============================================================
\section{Gravity Sensitivity Analysis}
\label{sec:gravity-sensitivity}
%% ============================================================

The critical insight motivating this paper is that the metal processing chain comprises distinct unit operations with fundamentally different gravity dependencies. We analyze each operation's gravity sensitivity to determine the minimum gravity level at which it becomes tractable.

\subsection{Slag--metal separation}

In terrestrial smelting, slag floats on molten metal due to density differences under Earth gravity ($\Delta\rho \approx 3{,}000$~kg/m$^3$ for iron--slag systems). The Stokes rise velocity of a slag droplet scales linearly with gravity:
\begin{equation}
    v_s = \frac{2 r^2 \Delta\rho\, g}{9 \mu}
    \label{eq:stokes}
\end{equation}
where $r$ is the droplet radius, $\Delta\rho$ the density difference, $g$ the gravitational acceleration, and $\mu$ the melt viscosity.

At Earth gravity, mm-scale slag droplets rise at velocities of order 1--10~cm/s, enabling continuous separation. In microgravity ($g \to 0$), this mechanism vanishes entirely. Alternative approaches---electromagnetic separation, centrifugal separation within the melt vessel, or acoustic manipulation---have been proposed but not demonstrated at scales relevant to industrial processing.

Critically, even very low gravity restores buoyancy separation. At $0.01g$, the Stokes velocity for a 1~mm slag droplet in liquid iron is approximately 0.1--1~mm/s, sufficient for separation over minutes rather than seconds. At $0.05g$, separation timescales approach terrestrial values for typical batch sizes.

\subsection{Convective heat transfer}

Terrestrial furnaces rely on natural convection to distribute heat throughout the melt and to establish the temperature gradients that drive solidification microstructure. The Grashof number, which characterizes the ratio of buoyancy to viscous forces, scales linearly with $g$:
\begin{equation}
    \mathrm{Gr} = \frac{g \beta \Delta T L^3}{\nu^2}
    \label{eq:grashof}
\end{equation}

In microgravity, $\mathrm{Gr} \to 0$ and heat transfer becomes purely conductive (Marangoni convection driven by surface tension gradients provides some mixing near free surfaces, but is insufficient for homogenization of bulk melts). At $0.01g$, the Grashof number for a 0.5~m melt pool at $\Delta T = 100$~K recovers to approximately $10^6$, well into the regime where natural convection provides meaningful mixing.

\subsection{Melt containment}

In microgravity, molten metal forms free-floating spheres that can be contained by electromagnetic, acoustic, or electrostatic forces. However, the power required for electromagnetic containment of a sphere of mass $m$ against perturbations scales as $m^{2/3}$ for surface confinement forces. For a 1~tonne melt pool, the required electromagnetic coil power becomes prohibitive (estimated 1--10~MW continuous), and the interaction between the EM field and the conductive melt introduces complex stirring patterns that may be detrimental to process control.

With even minimal gravity ($> 0.01g$), conventional refractory crucibles become viable. The gravity serves the same containment function as on Earth, albeit with modified buoyancy and convection patterns.

\subsection{Zone refining: a microgravity advantage}

Zone refining stands in contrast to the operations above. In the float zone process, a narrow molten zone traverses a solid rod, carrying impurities preferentially in the direction of travel according to their segregation coefficients. In terrestrial processing, buoyancy-driven convection within the molten zone partially remixes impurities, reducing purification efficiency. Suppressing this convection in microgravity should theoretically produce sharper impurity segregation, achieving solar-grade purity in fewer passes.

This advantage has been predicted theoretically and supported by analog experiments, but never demonstrated for silicon zone refining in sustained microgravity. The only space-based float zone silicon experiment (Dold \& Benz, sounding rocket, $\sim$6 minutes) was too brief for meaningful purification.

The UMG-Si pathway \cite{fornies2021} suggests that 4N--5N purity may be sufficient for viable solar cells (15--20\% efficiency), potentially reducing the number of zone refining passes needed and making the microgravity approach more practical.

\subsection{Summary of gravity sensitivity}

Table~\ref{tab:gravity-sensitivity} presents the gravity sensitivity analysis for each major unit operation.

\begin{table}[ht]
\centering
\caption{Gravity sensitivity by unit operation. ``Minimum viable $g$'' indicates the approximate threshold above which the operation becomes tractable using modified terrestrial approaches.}
\label{tab:gravity-sensitivity}
\begin{tabular}{lllp{4.5cm}}
\toprule
\textbf{Operation} & \textbf{Gravity need} & \textbf{Min.\ viable $g$} & \textbf{Rationale} \\
\midrule
Smelting & Required & 0.01--0.05$g$ & Buoyancy-driven slag separation \\
Slag separation & Required & 0.05--0.15$g$ & Stokes velocity $>$ 0.5~mm/s \\
Casting & Beneficial & 0.01$g$ & Mold filling, shrinkage control \\
Zone refining & Harmful & 0$g$ preferred & Convection degrades segregation \\
Thin-film CVD & Neutral & 0$g$ acceptable & Gas-phase transport dominates \\
Electrolysis & Magnetic alt.\ & 0$g$ acceptable & Diamagnetic separation \cite{akay2025} \\
\bottomrule
\end{tabular}
\end{table}

%% ============================================================
\section{Hybrid Multi-Gravity-Zone Architecture}
\label{sec:architecture}
%% ============================================================

The gravity sensitivity analysis in \S\ref{sec:gravity-sensitivity} motivates a station architecture with multiple gravity zones, matching the gravitational environment to the physics of each process step. We evaluate three candidate architectures.

\subsection{Architecture A: Pure microgravity}

All processing occurs in zero-$g$. Smelting uses electromagnetic levitation; slag separation uses EM or acoustic manipulation; zone refining exploits the convection-free environment.

\textbf{Advantages:} No rotating structures; simplest mechanical design; best environment for zone refining.

\textbf{Disadvantages:} Electromagnetic containment power for industrial-scale melts is prohibitive ($>$1~MW continuous for 1-tonne batches); no demonstrated pathway for slag separation at scale; thermal management without convection requires active forced circulation.

\textbf{Mass estimate:} 350,000--450,000~kg (large EM coil systems, high power requirement).

\textbf{Assessment:} Not viable at industrial scale. Physics limitations, not merely engineering gaps.

\subsection{Architecture B: Full artificial gravity}

A rotating station provides 0.3--0.5$g$ throughout, enabling near-terrestrial metallurgical processes with parametric adjustments.

\textbf{Advantages:} Directly applies terrestrial metallurgical knowledge; conventional equipment designs; well-understood convection and separation physics.

\textbf{Disadvantages:} Loses all microgravity advantages for refining; requires spinning a 400,000~kg station at significant RPM; Coriolis effects complicate fluid handling; incompatible with other station functions (docking, solar array pointing, antenna operations).

\textbf{Mass estimate:} 450,000--550,000~kg (structural penalty for centrifugal loads throughout station).

\textbf{Assessment:} Viable but wasteful. Imposes gravity where it is not needed and loses genuine microgravity advantages for refining.

\subsection{Architecture C: Hybrid multi-gravity-zone (recommended)}

A non-rotating core station with a dedicated rotating smelting arm. The architecture comprises:

\begin{enumerate}
    \item \textbf{Zero-$g$ core:} Central station with power generation, habitat, storage, docking. Also houses zone refining modules, thin-film deposition, and electrolysis (using magnetic separation).
    \item \textbf{Rotating smelting module:} A counter-rotating arm at $\sim$50~m radius, spinning at $\sim$1.4~RPM to produce 0.1$g$ at the smelting floor. Contains arc furnaces, slag separation vessels, and casting equipment.
    \item \textbf{Material transfer:} A central axle elevator moves feedstock outward to the smelting module and returns refined metal inward to the zero-$g$ core.
\end{enumerate}

\textbf{Advantages:} Optimal gravity environment for each operation; no need to spin the entire station; counter-rotation cancels angular momentum; smelting module is mechanically isolated from sensitive zero-$g$ operations.

\textbf{Disadvantages:} Rotating joint is a maintenance concern; Coriolis forces at 1.4~RPM are non-negligible for large melt pools; material transfer across the rotation interface requires engineering development.

\textbf{Mass estimate:} 340,000--430,000~kg. The rotating arm adds $\sim$25,000~kg of structural mass but eliminates $\sim$10,000--15,000~kg of electromagnetic confinement hardware that Architecture A would require, yielding a net penalty of only 10,000--15,000~kg.

\textbf{Assessment:} Recommended baseline. Best performance at lowest mass and risk.

\subsection{Architecture comparison}

\begin{table}[ht]
\centering
\caption{Architecture comparison for the Material Processing Station.}
\label{tab:architecture-comparison}
\begin{tabular}{lccc}
\toprule
\textbf{Metric} & \textbf{A: Pure 0$g$} & \textbf{B: Full AG} & \textbf{C: Hybrid} \\
\midrule
Station mass (tonnes) & 350--450 & 450--550 & 340--430 \\
Smelting viability & No & Yes & Yes \\
Zone refining quality & Best & Degraded & Best \\
EM containment power & $>$1~MW & None & None \\
Coriolis effects & None & Station-wide & Localized \\
Docking compatibility & Full & Constrained & Full \\
Technology readiness & TRL 2 & TRL 4--5 & TRL 3--4 \\
\bottomrule
\end{tabular}
\end{table}

%% ============================================================
\section{Gate 1 Decision Criteria}
\label{sec:gate1}
%% ============================================================

We propose quantitative go/no-go criteria for the Gate~1 technology decision at month~36 of the research program. These criteria are designed to be objectively measurable from the experimental campaigns described in \S\ref{sec:roadmap}.

\subsection{Mandatory criteria (all must pass)}

\begin{enumerate}
    \item \textbf{Slag separation efficiency:} $>$95\% metal recovery from a representative asteroid feedstock (CI/CM chondrite simulant) in $\leq$0.15$g$ at batch sizes $\geq$1~kg, demonstrated in at least 10 consecutive batches.
    \item \textbf{Zone refining purity:} Silicon zone-refined in microgravity achieves $\leq$1~ppmw total transition metal impurities (Fe+Cr+Cu) from MG-Si feedstock ($\sim$2\% purity) in $\leq$10 passes, demonstrated at $\geq$100~g batch size.
    \item \textbf{Electrolysis efficiency:} Water electrolysis with magnetic separation achieves $\geq$80\% of terrestrial current density at $\geq$1~kW scale in sustained ($>$100~hr) microgravity operation.
    \item \textbf{Structural integrity:} Rotating smelting arm structural and bearing design passes preliminary design review (PDR) with positive margins for 10-year operational life at 1.4~RPM.
\end{enumerate}

\subsection{Desirable criteria (inform architecture optimization)}

\begin{enumerate}
    \item Minimum viable gravity for slag separation determined to $\pm$0.02$g$ precision.
    \item Coriolis effects on melt pool homogeneity characterized at 1--3~RPM for $\geq$10~kg melt pools.
    \item Autonomous process control demonstrated for $\geq$48~hr continuous smelting without human intervention.
\end{enumerate}

%% ============================================================
\section{Research Roadmap}
\label{sec:roadmap}
%% ============================================================

We propose a phased research program to retire the critical unknowns:

\textbf{Phase~1 (Years 1--2, \$100--160M):} Establish a terrestrial analog facility using electromagnetic levitation and representative asteroid simulant compositions. Conduct hundreds of smelting, separation, and refining trials to optimize slag chemistry, flux formulations, and autonomous control algorithms. Commission CFD modeling calibrated to terrestrial data.

\textbf{Phase~2 (Years 2--4, \$300--400M):} Design and deploy a dedicated partial-gravity experiment platform---potentially a free-flying centrifuge rather than ISS-hosted---capable of sustaining 0.01--0.2$g$ for hours to days. Conduct smelting and slag separation experiments at 1--10~kg batch scale across a systematic gravity sweep. Simultaneously conduct ISS microgravity zone refining experiments at 100~g to 1~kg scale.

\textbf{Phase~3 (Years 4--6, \$150--250M):} Based on Phase~2 results, build a rotating smelting module engineering development unit for ground testing. Conduct integrated processing chain demonstrations: feedstock preparation $\to$ smelting $\to$ slag separation $\to$ zone refining $\to$ product characterization. Validate Gate~1 criteria.

Total estimated cost: \$550--810M over 6 years. This is approximately 1--1.6\% of the Material Processing Station's estimated \$50B capital cost.

%% ============================================================
\section{Discussion}
\label{sec:discussion}
%% ============================================================

\subsection{The reframing from risk to engineering}

The most significant outcome of this analysis is the reframing of microgravity metallurgy from a project-ending research risk (TRL 2--3) to an engineering development challenge (TRL 3--4 with clear maturation pathway). This reframing hinges entirely on the hybrid architecture: by accepting that bulk smelting requires gravity and providing it locally through a rotating module, the problem decomposes into a set of challenges that are individually tractable, even if collectively demanding.

The contingency allocation for the processing system can correspondingly be reduced from the 30--40\% appropriate for a research-dominated program to the 15--20\% appropriate for engineering development, freeing approximately \$5--10B in the overall program budget.

\subsection{Implications for station design}

The hybrid architecture has several downstream implications:

\begin{enumerate}
    \item \textbf{Power allocation:} Elimination of multi-MW electromagnetic containment systems reduces the processing power budget by an estimated 40--60\%, freeing capacity for production scale-up.
    \item \textbf{Thermal management:} Convective heat transfer in the rotating module simplifies thermal design relative to the pure microgravity case, where all heat transfer must be actively managed.
    \item \textbf{Feedstock flexibility:} Gravity-based processing is more tolerant of feedstock variability than electromagnetic processing, which requires tight control of electrical conductivity and magnetic susceptibility.
\end{enumerate}

\subsection{Limitations}

This analysis has several limitations. First, our gravity sensitivity thresholds are derived from dimensional analysis and extrapolation from 1$g$ data; they have not been validated experimentally in the 0.01--0.2$g$ regime. Second, the mass estimates for the three architectures carry uncertainty of approximately $\pm$30\%. Third, we have not addressed the thermal design of the rotating module in detail, particularly the challenge of dissipating waste heat from a continuously operating smelting facility through a rotating joint. Fourth, the Coriolis effects on large melt pools at 1--3~RPM are genuinely unknown and could introduce process complications not captured by our analysis.

%% ============================================================
\section{Conclusion}
\label{sec:conclusion}
%% ============================================================

The question ``Can microgravity metallurgy scale to industrial production?'' has a nuanced answer: not in its current form, but the challenge can be resolved through architectural design rather than fundamental research breakthroughs. By decomposing the metal processing chain into gravity-sensitive and gravity-independent operations and providing the appropriate gravitational environment for each, a hybrid multi-gravity-zone station achieves industrial-scale processing within the existing mass and cost envelope.

The critical path to realizing this architecture is experimental characterization of metallurgical processes in the 0.01--0.2$g$ regime---a gravity range with virtually no existing data. We have proposed a \$550--810M, 6-year research program to fill this gap and defined quantitative Gate~1 criteria that can objectively determine whether to proceed with station construction.

The shift from pure microgravity to hybrid gravity processing represents a pragmatic compromise: it sacrifices the elegance of a fully zero-$g$ factory for the reliability of proven physical principles adapted to a new environment. This is, we argue, exactly the right trade-off for a program that cannot afford a project-ending technology failure.

%% ============================================================
%% References
%% ============================================================

\begin{thebibliography}{99}

\bibitem{hakkenberg2026consensus}
T.\ Hakkenberg, ``Multi-Model AI Consensus for Engineering Decision-Making: Methodology and Validation,'' \textit{in preparation}, 2026.

\bibitem{ratke1995}
L.\ Ratke and S.\ Diefenbach, ``Liquid immiscible alloys,'' \textit{Materials Science and Engineering: R: Reports}, vol.~15, no.~7--8, pp.~263--347, 1995.

\bibitem{dangelo2021}
O.\ D'Angelo, F.\ Kuthe, and S.-J.\ Liu, ``A gravity-independent powder-based additive manufacturing process tailored for space applications,'' arXiv:2102.09815, 2021.

\bibitem{todaro2020}
C.J.\ Todaro, M.A.\ Easton, D.\ Qiu et al., ``Grain refinement of stainless steel in ultrasound-assisted additive manufacturing,'' arXiv:2008.04485, 2020.

\bibitem{akay2025}
\"O.\ Akay, M.\ Monfort-Castillo, T.\ St~Francis et al., ``Magnetically induced convection enhances water electrolysis in microgravity,'' \textit{Nature Chemistry}, vol.~17, pp.~1673--1679, 2025.

\bibitem{romerocalo2022}
\'A.\ Romero-Calvo, \"O.\ Akay, H.\ Schaub, and K.\ Brinkert, ``Magnetic phase separation in microgravity,'' \textit{npj Microgravity}, vol.~8, pp.~1--10, 2022.

\bibitem{fornies2021}
E.\ Forni\'es, C.\ del Ca\~nizo, L.\ M\'endez et al., ``UMG silicon for solar PV: from defects detection to PV module degradation,'' arXiv:2101.08019, 2021.

\bibitem{freitas1982}
R.A.\ Freitas Jr.\ and W.P.\ Gilbreath (eds.), ``Advanced Automation for Space Missions,'' NASA Conference Publication CP-2255, 1982.

\bibitem{metzger2013}
P.T.\ Metzger, A.\ Muscatello, R.P.\ Mueller, and J.\ Mantovani, ``Affordable, Rapid Bootstrapping of the Space Industry and Solar System Civilization,'' \textit{Journal of Aerospace Engineering}, vol.~26, no.~1, pp.~18--29, 2013.

\bibitem{hofer2019}
R.R.\ Hofer, H.\ Kamhawi, I.G.\ Mikellides et al., ``Completing the Development of the 12.5~kW Hall Effect Rocket with Magnetic Shielding,'' IEPC-2019-193, 36th International Electric Propulsion Conference, 2019.

\bibitem{frieman2019}
J.D.\ Frieman, H.\ Kamhawi, J.A.\ Mackey et al., ``Completion of the Long Duration Wear Test of the NASA HERMeS Hall Thruster,'' AIAA-2019-3895, AIAA Propulsion and Energy Forum, 2019.

\bibitem{goebel2008}
D.M.\ Goebel and I.\ Katz, \textit{Fundamentals of Electric Propulsion: Ion and Hall Thrusters}, Wiley/JPL, 2008.

\end{thebibliography}

\end{document}
